% GENERATED FILE. Edit canonical lessons, then run npm run generate:books.
% canonical-source-hash: fnv1a64:87c07d880bf1fb2b
% canonical-lessons: ZH-C304-ruxue, ZH-C304-riji, ZH-W304-ji-note, ZH-C304-baodao-register, ZH-C304-xuefei, ZH-C304-jizhu

\chapter{Starting School}
\label{ch:zh-starting-school}
\hlchaptermodality{304}

\begin{chapteropening}
\textbf{By the end of this chapter:} \emph{I can talk about starting school, signing in, the fees, and keeping a diary.}

Complete the last lesson of chapter 304: I can talk about starting school, signing in, the fees, and keeping a diary.
\end{chapteropening}

\section[rùxué]{rùxué — to start school}

\label{lesson:ZH-C304-ruxue}



\begin{quote}
Before the new one: say the Mandarin for to start the school term, then the Mandarin for a new student at school.
\end{quote}

\subsection*{You'll want to know: rùxué}
\emph{rùxué} — \textquotedblleft{}to start school\textquotedblright{}.

\begin{quote}
\emph{tā liù suì rùxué} — He started school at six.
\end{quote}

Its characters come later; for now it is a word for the ear.

\begin{cousinweb}
\emph{rùxué} is built from \emph{rù} 'to enter' and \zh{学} 'to study', both inherited from Middle Chinese, \emph{nyip} and \emph{haewk} (Guangyun, in Baxter's transcription), which glosses the first 'inside' and 'to take in'.
\end{cousinweb}

\begin{grammarlens}[title={What you've built}]
One more word for telling what happened.
\end{grammarlens}

\subsection*{Your turn}
\begin{itemize}
\raggedright
  \item \emph{Say it:} \emph{rùxué}
  \item \emph{Say it:} \emph{rùxué}, once more
  \item \emph{Say it:} \emph{tā liù suì rùxué} — He started school at six.
  \item \emph{Recall:} say the Mandarin for to start the school term, then the Mandarin for a new student at school, then say \emph{rùxué} again
\end{itemize}

\begin{tcolorbox}[breakable,colback=teal!4,colframe=teal!35!black,title={Before you move on}]
What does rùxué mean? (\textquotedblleft{}To start school\textquotedblright{}.) Say it once more.
\end{tcolorbox}
\section[\texorpdfstring{\zh{日记} (rìjì)}{日记 (rìjì)}]{\zh{日记} (rìjì) — diary}

\label{lesson:ZH-C304-riji}



\begin{quote}
Before the new one: say the Mandarin for a new student at school, then the Mandarin for to start school as a child.
\end{quote}

\subsection*{You'll want to know: \zh{日记}}
\textbf{\zh{日记}} — \emph{rìjì} — \textquotedblleft{}diary\textquotedblright{}.

\begin{quote}
\textbf{\zh{我小时候写过日记}} \emph{wǒ xiǎoshíhou xiě guo rìjì} — I kept a diary when I was little.
\end{quote}

The new character, in pieces: \zh{记} — five strokes in two pieces: the speech piece of \zh{说} on the left, two strokes; on the right the right piece of \zh{起}, three strokes. The next lesson puts it in your hand.

\begin{cousinweb}
\textbf{\zh{日记}} (\emph{rìjì}) is built from \zh{日} 'a day' and \zh{记} 'to note down', both inherited from Middle Chinese, \emph{nyit} and \emph{kiH} (Guangyun, in Baxter's transcription), which glosses the second 'to record'.
\end{cousinweb}

\begin{grammarlens}[title={What you've built}]
One more word for telling what happened.
\end{grammarlens}

\subsection*{Your turn}
\begin{itemize}
\raggedright
  \item \emph{Say it:} \emph{rìjì}
  \item \emph{Say it:} \emph{rìjì}, once more
  \item \emph{Say it:} \emph{wǒ xiǎoshíhou xiě guo rìjì} — I kept a diary when I was little.
  \item \emph{Recall:} say the Mandarin for a new student at school, then the Mandarin for to start school as a child, then say \emph{rìjì} again
\end{itemize}

\begin{tcolorbox}[breakable,colback=teal!4,colframe=teal!35!black,title={Before you move on}]
What does \zh{日记} mean? (\textquotedblleft{}Diary\textquotedblright{}.) Say it once more.
\end{tcolorbox}
\section[\texorpdfstring{\zh{记} (jì)}{记 (jì)}]{\zh{记} — the speech piece, then the right piece of \zh{起}}

\label{lesson:ZH-W304-ji-note}



\begin{quote}
The word first, then the character inside it.

\begin{itemize}
\raggedright
  \item \emph{Recall:} say \emph{rìjì} — \textbf{R1}, one lesson back, before it decays
\end{itemize}

\begin{quote}
\zh{说} \zh{起}
\end{quote}

\emph{Read it:} each one aloud, with nothing said first — \emph{shuō}, \emph{qǐ}
\end{quote}

\subsection*{Script — the speech piece, then the right piece of \zh{起}}
\hlblockfigure{\detokenize{figures/ZH-W304-ji-note-filmstrip.pdf}}{How \zh{记} is written, stroke by stroke}

\begin{quote}
\zh{记}
\end{quote}

\textbf{Two pieces, left before right.} Count the pieces, not the strokes.

\textbf{The left piece, two strokes.} The speech piece of \zh{说}: a dot, then a short stroke that runs right, turns down, and rises to the right at the foot.

\textbf{The right piece, three strokes}, as on the right of \zh{起}: a bar that turns down, a bar across the middle, then a vertical that bends right along the foot and hooks up.

2 + 3 = five.

\subsection*{Your turn}
\begin{itemize}
\raggedright
  \item \emph{Write it:} \zh{记}, all five, with the model in front of you
  \item \emph{Cover:} write it again from memory, and count the pieces, not the strokes
\end{itemize}

\begin{tcolorbox}[breakable,colback=teal!4,colframe=teal!35!black,title={Before you move on}]
How many pieces in \zh{记}? (\textbf{Two}.) What is the left piece? (\textbf{The speech piece, as in \zh{说}}.) Where have you made the right piece? (\textbf{On the right of \zh{起}}.)

Source: \href{https://raw.githubusercontent.com/chanind/hanzi-writer-data/68d10a4b21150cae5e1ebbd223eed289cf32d90c/data/\%E8\%AE\%B0.json}{Hanzi Writer Data \zh{记}, pinned stroke paths and medians}.
\end{tcolorbox}
\section[\texorpdfstring{\zh{报到} (bàodào)}{报到 (bàodào)}]{\zh{报到} (bàodào) — to report in on arrival; to sign in}

\label{lesson:ZH-C304-baodao-register}



\begin{quote}
Before the new one: say the Mandarin for to start school as a child, then the Mandarin for a diary.
\end{quote}

\subsection*{You'll want to know: \zh{报到}}
\textbf{\zh{报到}} — \emph{bàodào} — \textquotedblleft{}to report in on arrival; to sign in\textquotedblright{}.

\begin{quote}
\textbf{\zh{新生明天来学校报到}} \emph{xīnshēng míngtiān lái xuéxiào bàodào} — The new students sign in at school tomorrow.
\end{quote}

It sounds the same as \emph{bàodào} 'to report the news'; this one means arriving and signing in.

\begin{cousinweb}
\textbf{\zh{报到}} (\emph{bàodào}) is built from \zh{报} 'to report' and \zh{到} 'to arrive', both inherited from Middle Chinese, \emph{pawH} and \emph{tawH} (Guangyun, in Baxter's transcription), which glosses the second 'to reach': to report that you have arrived.
\end{cousinweb}

\begin{grammarlens}[title={What you've built}]
One more word for telling what happened.
\end{grammarlens}

\subsection*{Your turn}
\begin{itemize}
\raggedright
  \item \emph{Say it:} \emph{bàodào}
  \item \emph{Say it:} \emph{bàodào}, once more
  \item \emph{Say it:} \emph{xīnshēng míngtiān lái xuéxiào bàodào} — The new students sign in at school tomorrow.
  \item \emph{Recall:} say the Mandarin for to start school as a child, then the Mandarin for a diary, then say \emph{bàodào} again
\end{itemize}

\begin{tcolorbox}[breakable,colback=teal!4,colframe=teal!35!black,title={Before you move on}]
What does \zh{报到} mean? (\textquotedblleft{}To report in on arrival; to sign in\textquotedblright{}.) Say it once more.
\end{tcolorbox}
\section[xuéfèi]{xuéfèi — tuition fees}

\label{lesson:ZH-C304-xuefei}



\begin{quote}
Before the new one: say the Mandarin for a diary, then the Mandarin for to report in on arrival.
\end{quote}

\subsection*{You'll want to know: xuéfèi}
\emph{xuéfèi} — \textquotedblleft{}tuition fees\textquotedblright{}.

\begin{quote}
\emph{dàxué de xuéfèi hěn guì} — University fees are high.
\end{quote}

Its characters come later; for now it is a word for the ear.

\begin{cousinweb}
\emph{xuéfèi} is built from \zh{学} 'to study' and \emph{fèi} 'to spend', both inherited from Middle Chinese, \emph{haewk} and \emph{phj+jH} (Guangyun, in Baxter's transcription), which glosses the second 'to use up'.
\end{cousinweb}

\begin{grammarlens}[title={What you've built}]
One more word for telling what happened.
\end{grammarlens}

\subsection*{Your turn}
\begin{itemize}
\raggedright
  \item \emph{Say it:} \emph{xuéfèi}
  \item \emph{Say it:} \emph{xuéfèi}, once more
  \item \emph{Say it:} \emph{dàxué de xuéfèi hěn guì} — University fees are high.
  \item \emph{Recall:} say the Mandarin for a diary, then the Mandarin for to report in on arrival, then say \emph{xuéfèi} again
\end{itemize}

\begin{tcolorbox}[breakable,colback=teal!4,colframe=teal!35!black,title={Before you move on}]
What does xuéfèi mean? (\textquotedblleft{}Tuition fees\textquotedblright{}.) Say it once more.
\end{tcolorbox}
\section[jìzhù]{jìzhù — to commit to memory; to keep in mind}

\label{lesson:ZH-C304-jizhu}



\begin{quote}
Before the new one: say the Mandarin for to report in on arrival, then the Mandarin for school tuition fees.
\end{quote}

\subsection*{You'll want to know: jìzhù}
\emph{jìzhù} — \textquotedblleft{}to commit to memory; to keep in mind\textquotedblright{}.

\begin{quote}
\emph{wǒ jìzhù le tā de diànhuà hàomǎ} — I memorized his phone number.
\end{quote}

Its characters come later; for now it is a word for the ear.

\begin{cousinweb}
\emph{jìzhù} is built from \zh{记} 'to note down' and \emph{zhù} 'to stay', both inherited from Middle Chinese, \emph{kiH} and \emph{drjuH} (Guangyun, in Baxter's transcription), which glosses the second 'to stop': noted so that it stays.
\end{cousinweb}

\begin{grammarlens}[title={What you've built}]
One more word for telling what happened.
\end{grammarlens}

\subsection*{Your turn}
\begin{itemize}
\raggedright
  \item \emph{Say it:} \emph{jìzhù}
  \item \emph{Say it:} \emph{jìzhù}, once more
  \item \emph{Say it:} \emph{wǒ jìzhù le tā de diànhuà hàomǎ} — I memorized his phone number.
  \item \emph{Recall:} say the Mandarin for to report in on arrival, then the Mandarin for school tuition fees, then say \emph{jìzhù} again
\end{itemize}

\begin{tcolorbox}[breakable,colback=teal!4,colframe=teal!35!black,title={Before you move on}]
What does jìzhù mean? (\textquotedblleft{}To commit to memory; to keep in mind\textquotedblright{}.) Say it once more.
\end{tcolorbox}
